\section{A finite tag clock}
\label{sec:clock}

The affine-line lemma requires the two factors to communicate finite data.
We call a possible collection of these data a \emph{tag}. This section
constructs a clock that controls independent tag choices by testing the
two length residues, measured from scheduled reference positions. Its relation permits either matching
tags or a single specified mismatch. The latter can later be absorbed into
an affine correction. A second property bounds ambiguous pairs of candidate
origins uniformly in their common translation. That bound must depend only
on the tags, since the number of candidate origins will be chosen afterward.

\begin{lemma}[Finite tag clock]\label{lem:clock}
Let $\Lambda$ be a finite nonempty set, let $\mathcal R$ be any finite set,
and let $Q\geq 1$ be an integer.  Put
\[
 F_{\rm cl}=400,
 \qquad
 C_{\rm cl}=2^{|\Lambda|}\bigl(1+4F_{\rm cl}^{2}\bigr).
\]
There exist a positive integer $n$, all of whose prime factors exceed $Q$,
sets $I_\alpha,J_\alpha\subseteq\mathbb Z/n\mathbb Z$ for
$\alpha\in\Lambda$, a set $W\subseteq\mathbb Z/n\mathbb Z$, and residues
$v_p\in\mathbb Z/n\mathbb Z$ for $p\in\mathcal R$ with the following
properties.
\begin{enumerate}[label=\textup{(\roman*)}]
\item For every $s\in\mathbb Z/n\mathbb Z$, the relation
\[
 \mathcal E(s)=
 \{(\alpha,\alpha')\in\Lambda^2:
             s\in I_\alpha+J_{\alpha'}\}
\]
is either a subset of the diagonal or a singleton consisting of an
off-diagonal pair.  If $s\notin W$, it is the full diagonal.
\item For every $u\in\mathbb Z/n\mathbb Z$,
\[
 \#\{(p,p')\in\mathcal R^2:p\ne p',\quad
                 u+v_{p'}-v_p\in W\}\leq C_{\rm cl}.
\]
\end{enumerate}
In particular, $C_{\rm cl}$ depends only on $|\Lambda|$.
\end{lemma}

\begin{proof}
The construction uses one block of coordinates for every left/right
assignment of the tags. First-coordinate restrictions exclude competing
tag pairs, while a batch of additional coordinates excludes reverse pairs
and controls translated roster differences. We prove these two tasks
separately and only then identify the product group with one cyclic group.

\emph{Coordinate blocks and their sets.}
We first construct the sets in a finite product of cyclic groups.  For each
map $\chi:\Lambda\to\{\mathrm L,\mathrm R\}$, reserve one first
coordinate and a batch of coordinates.  The batch contains a dedicated
coordinate for every ordered list
\[
 \bigl((p_a,p'_a)\bigr)_{a=0}^{F_{\rm cl}-1}
\]
of ordered pairs of distinct elements of $\mathcal R$ such that the
underlying unordered edges are distinct and form a forest.  Add a spare
coordinate if necessary so that the batch has positive odd size.  Thus, even
when no such list exists, the batch consists of one spare coordinate.
There are finitely many coordinates.  Give them distinct prime orders $q$,
each satisfying
\[
 q>\max\{Q,400,2^{|\mathcal R|+1}\}.
\]
There are enough primes: the product-plus-one argument supplies a prime
outside any prescribed finite list.  Write $G$ for the resulting product of
cyclic groups.

Identify a coordinate of order $q$ with the subgroup
$q^{-1}\mathbb Z/\mathbb Z$ of the circle
$\mathbb T=\mathbb R/\mathbb Z$.  Write $\|x\|_{\mathbb T}$ for circular
distance to zero, and call a coordinate value \emph{small} if
$\|x\|_{\mathbb T}\leq 1/100$.  Define the two sets of centers
\[
 C_I=\{0,1,3\}/6,
 \qquad C_J=\{2,4,5\}/6
 \quad\text{in }\mathbb T.
\]
For each tag $\alpha$, impose the following constraints in the block
belonging to $\chi$.
\begin{itemize}
\item If $\chi(\alpha)=\mathrm L$, require the first coordinate of
  $I_\alpha$ to be zero and that of $J_\alpha$ to be nonzero.  In every
  batch coordinate, require the $I_\alpha$ value to have distance at most
  $1/16$ from $C_I$, and the $J_\alpha$ value to have distance at most
  $1/16$ from $C_J$.
\item If $\chi(\alpha)=\mathrm R$, require the first coordinate of
  $I_\alpha$ to be nonzero and that of $J_\alpha$ to be zero.  For each of
  $I_\alpha$ and $J_\alpha$, require a strict majority of the batch
  coordinates to be small.
\end{itemize}
The constraints over all blocks define $I_\alpha,J_\alpha\subseteq G$.
Let $W\subseteq G$ be the union, over all $\chi$, of the following two
obstructions: the first coordinate is zero; or no coordinate of the batch
is small.

\emph{Sums and differences inside one batch.}
We verify the sum and difference properties of a batch.  For a left tag,
every target in each cyclic coordinate is a sum of allowed values.  Indeed,
$C_I+C_J$ contains all six multiples of $1/6$.  Given a target $t$, choose
$c_I\in C_I$, $c_J\in C_J$ and a signed displacement $\delta$ with
$|\delta|\leq 1/12$ such that
$t=c_I+c_J+\delta$ in $\mathbb T$.  Round
$c_I+\delta/2$ to a nearest grid point $i$, and put $j=t-i$.  Both $i$ and
$j$ are grid points, and their distances from $c_I$ and $c_J$, respectively,
are at most
\[
 \frac1{24}+\frac1{2q}
 <\frac1{24}+\frac1{800}<\frac1{16}.
\]
Moreover, the distance between any center in $C_I$ and any center in $C_J$
is at least $1/6$.  Consequently every allowed left-tag difference $i-j$
satisfies, in every batch coordinate,
\[
 \|i-j\|_{\mathbb T}\geq
 \frac16-\frac2{16}=\frac1{24}.
\]

For a right tag, any target batch with one small coordinate is a sum of
allowed batches.  To see this, write the batch size as $2k+1$ and choose a
small target coordinate.  At that coordinate, assign the target to the
first summand and zero to the second.  Partition the other $2k$ coordinates
into two sets of size $k$.  On the first set put zero in the first summand
and the target in the second; on the second set do the reverse.  Both
summands then have at least $k+1$ small coordinates.  Conversely, any two
allowed right-tag batches have a common small coordinate, because their
sets of small coordinates are strict majorities.  Their difference
therefore has some coordinate satisfying
\[
 \|i-j\|_{\mathbb T}\leq\frac2{100}<\frac1{24}.
\]
Thus the difference sets $I_\alpha-J_\alpha$ and
$I_{\alpha'}-J_{\alpha'}$ are disjoint whenever a block places $\alpha$
on the left and $\alpha'$ on the right.

\emph{The full relation, including independently mismatched tags.}
If $s\notin W$, all first coordinates of $s$ are nonzero and every batch
has a small coordinate.  The preceding sum constructions, together with
the first-coordinate decompositions $s=0+s$ and $s=s+0$, show that
$s\in I_\alpha+J_\alpha$ for every tag $\alpha$.  The choices in distinct
blocks are independent, so they combine to give elements of $G$.

Now suppose an off-diagonal pair $(\alpha,\alpha')$ belongs to
$\mathcal E(s)$.  Consider any competing pair $(\lambda,\mu)$ with
$\{\lambda,\mu\}\ne\{\alpha,\alpha'\}$.  There is an assignment
$\chi$ placing $\alpha$ on the left, $\alpha'$ on the right, and
$\lambda,\mu$ on the same side: identifying $\lambda$ with $\mu$ does
not identify $\alpha$ with $\alpha'$.  In its first coordinate the pair
$(\alpha,\alpha')$ forces $s=0$, whereas the pair $(\lambda,\mu)$
requires $s\ne0$.  Hence that competitor is impossible.  The only other
possible competitor is the reverse pair $(\alpha',\alpha)$.  If both
pairs represented $s$, there would be
\[
 i_\alpha+j_{\alpha'}=i_{\alpha'}+j_\alpha,
 \qquad\text{hence}\qquad
 i_\alpha-j_\alpha=i_{\alpha'}-j_{\alpha'},
\]
with the indicated elements belonging to their tag sets.  A block
separating these tags contradicts the disjointness of the difference
sets.  This proves the first alternative in \textup{(i)}.  It also excludes
off-diagonal pairs outside $W$, where the full diagonal is present, and
completes the proof of \textup{(i)}.

\emph{Assigning the roster and counting whole-batch obstructions.}
The tag relation is now established. What remains is a choice of the
$v_p$ that makes membership in $W$ rare among ordered roster differences,
for every common translate.  Write
$r=|\mathcal R|$ and enumerate the roster as $p_0,\ldots,p_{r-1}$.  In
every first coordinate set the residue of $v_{p_j}$ equal to $2^j$.  All
ordered differences for distinct roster slots are distinct modulo that
coordinate's prime.  For $r\geq2$, they are distinct as integers: the sign
determines their orientation, and in a positive difference
\[
 2^b-2^a=2^a(2^{b-a}-1),\qquad a<b,
\]
the power of two dividing it determines $a$, after which the odd factor
determines $b$.  Reduction modulo $q$ creates no collision, since the
distance between any two such integer differences is less than $2^r<q$.
The assertion is vacuous when $r<2$.  For any translate $u$, at most one
ordered pair therefore hits a specified first-coordinate zero obstruction.

In the dedicated batch coordinate for a list
$((p_a,p'_a))_{a=0}^{F_{\rm cl}-1}$, let $z_a$ be a grid point at circular
distance at most $1/(2q)$ from $a/F_{\rm cl}$.  Prescribe
\[
 (v_{p'_a}-v_{p_a})\big|_{\text{coordinate}}=z_a
 \qquad(0\leq a<F_{\rm cl}).
\]
These prescriptions can be satisfied: choose a root value in each tree
and assign successive vertex values along its edges, with the sign
determined by the prescribed orientation.  A forest has no consistency
condition around a cycle.  Assign arbitrary values to any remaining
vertices and in spare coordinates.  All these assignments are made
independently in each coordinate, and so define elements $v_p\in G$.

For any coordinate translate $t$, the equally spaced points
$a/F_{\rm cl}$ include one at distance at most $1/(2F_{\rm cl})$ from
$-t$.  The corresponding prescribed difference satisfies
\[
 \|t+z_a\|_{\mathbb T}
 \leq\frac1{2F_{\rm cl}}+\frac1{2q}
 <\frac1{400}<\frac1{100}.
\]
Thus every dedicated list has a pair whose translated difference is small
in its dedicated coordinate, uniformly in the translate.

Fix $u\in G$ and a batch.  Call an ordered pair $(p,p')$, $p\ne p'$,
bad for this batch if $u+v_{p'}-v_p$ has no small coordinate in the batch.
Suppose there were more than $4F_{\rm cl}^{2}$ such pairs.  Form the
simple undirected graph consisting of their underlying edges, discarding
isolated vertices.  If a spanning forest had fewer than $F_{\rm cl}$
edges, its nontrivial components would together have at most
$2F_{\rm cl}-2$ vertices: a component with $h\geq2$ vertices contributes
$h-1$ forest edges and satisfies $h\leq2(h-1)$.  The number of ordered
pairs would then be at most
\[
 (2F_{\rm cl}-2)(2F_{\rm cl}-3)<4F_{\rm cl}^{2},
\]
a contradiction.  Select $F_{\rm cl}$ edges of a spanning forest and
orient each as one of the available bad pairs.  In any order, they form a
list with a dedicated coordinate in this batch.  The preceding covering
estimate makes one of these pairs small in that coordinate, again a
contradiction.  Each batch therefore has at most $4F_{\rm cl}^{2}$ bad
ordered pairs.  This conclusion also applies when the roster is too small
to admit a dedicated list, by the same counting argument.

There are $2^{|\Lambda|}$ blocks.  Taking the union of their
first-coordinate and batch obstructions gives, for every $u\in G$,
\[
 \#\{(p,p'):p\ne p',\ u+v_{p'}-v_p\in W\}
 \leq 2^{|\Lambda|}\bigl(1+4F_{\rm cl}^{2}\bigr).
\]
For the empty roster the left side is zero throughout the construction.

\emph{One cyclic clock.}
Finally, let $n$ be the product of all coordinate primes.  The residue map
from $\mathbb Z/n\mathbb Z$ to $G$ is an injective homomorphism, since the
distinct primes are pairwise coprime, and both groups have $n$ elements.
It is therefore an isomorphism.  Transport all constructed sets and
residues through it.  This preserves sums, differences, and the two
properties, while every prime factor of $n$ exceeds $Q$.  All sets and
choices in the construction are finite.
\end{proof}
